\documentclass[11pt]{article}

% "final" genera la versión final sin marcas de revisión anónima.
\usepackage[final]{acl}

\usepackage{times}
\usepackage{latexsym}
\usepackage[T1]{fontenc}
\usepackage[utf8]{inputenc}
\usepackage{microtype}
\usepackage{inconsolata}
\usepackage{graphicx}
\usepackage{booktabs}
\usepackage{multirow}
\usepackage{amsmath}
\usepackage{caption}

\title{Evaluación de Robustez Léxica en Modelos de Lenguaje\\
  mediante Perturbaciones Controladas en Tareas de\\
  Razonamiento Verbal en Español}

\author{
  \textbf{Grupon 4?} \\
  - cmep
}

\begin{document}
\maketitle

% ─────────────────────────────────────────────────────────────────────────────
\begin{abstract}
Los modelos de lenguaje de gran escala exhiben rendimiento notable en tareas de
comprensión verbal, pero sus mecanismos de inferencia siguen siendo parcialmente
opacos: no es evidente si el modelo razona sobre el contenido semántico o si
explota patrones superficiales del texto.
En este trabajo presentamos un marco sistemático para evaluar la
\textbf{robustez léxica} de modelos de lenguaje en tareas de razonamiento verbal
en español.
Se implementan cuatro técnicas de perturbación controlada —sustitución por
sinónimos, parafraseo, pares mínimos y eliminación de atajos lingüísticos—
organizadas como módulos independientes que derivan de una clase abstracta común.
Las perturbaciones se aplican sobre un conjunto de datos de 13\,829 ejemplos en
español que abarca siete tipos de tarea.
La validez semántica de cada perturbación se verifica automáticamente mediante
similitud coseno de \textit{embeddings} multilingües, garantizando que el
significado y la etiqueta correcta se preserven.
El marco genera conjuntos de datos perturbados en tres niveles de intensidad
(bajo, medio, alto) para medir, en una fase posterior, el impacto sobre la
precisión y la tasa de cambio de predicciones de modelos evaluadores.
\end{abstract}

% ─────────────────────────────────────────────────────────────────────────────
\section{Introducción}
\label{sec:intro}


Los modelos de lenguaje de gran escala (LLMs) basados en la arquitectura
\textit{transformer} \citep{devlin-etal-2019-bert} han transformado el
procesamiento del lenguaje natural y, en pocos años, se han convertido en
componentes centrales de sistemas de búsqueda, asistencia educativa,
evaluación automática y toma de decisiones.
Su despliegue masivo, sin embargo, ha superado al ritmo al que la comunidad
puede caracterizar formalmente \emph{qué} es lo que estos modelos están
aprendiendo realmente: persiste una brecha entre la fluidez aparente de sus
respuestas y la certeza de que dicha fluidez proviene de una comprensión
semántica genuina.

Esta preocupación se hace especialmente relevante en tareas de
\textbf{razonamiento verbal}, dominio en el cual el modelo debe interpretar
relaciones lógicas, inferenciales y pragmáticas entre fragmentos de texto.
Pruebas de aptitud académica empleadas en contextos hispanohablantes
—comprensión lectora, ordenamiento y eliminación de oraciones, analogías,
series verbales y oraciones incompletas— constituyen un escenario natural
para examinar dicho razonamiento, ya que fueron diseñadas por humanos para
medir precisamente esa capacidad en personas.

El escenario, no obstante, está marcado por un desequilibrio importante:
mientras que la mayoría de los \textit{benchmarks} y herramientas de
evaluación de robustez se han desarrollado y validado en
inglés~\citep{jia-liang-2017-adversarial,jin-etal-2020-bert,ribeiro-etal-2020-checklist},
el español sigue subrepresentado en este tipo de análisis.
Esto deja un vacío metodológico significativo: los modelos que se utilizan
en aplicaciones hispanohablantes no son evaluados con el mismo nivel de
rigor adversarial que sus contrapartes en inglés, a pesar de operar sobre
hablantes y datos con características lingüísticas propias.

\section{Problema abordado}
\label{sec:problem}

Estudios recientes han demostrado que una fracción importante del rendimiento
de los LLMs en tareas de comprensión y razonamiento se apoya en
\textit{atajos léxicos}: correlaciones superficiales entre palabras del
enunciado y la respuesta correcta, en lugar de una representación profunda
del significado~\citep{mccoy-etal-2019-right,gardner-etal-2020-evaluating}.
Como consecuencia, modificaciones mínimas del texto(sustituir una palabra
por un sinónimo, parafrasear una oración, invertir un conector causal o
eliminar una pista discursiva) pueden alterar drásticamente la predicción
del modelo, aun cuando el significado y la respuesta correcta permanezcan
intactos.

Este fenómeno plantea un problema concreto: \textbf{no es posible distinguir,
únicamente a partir de la precisión global, si un modelo de lenguaje en
español razona sobre el contenido o si meramente explota regularidades
superficiales del texto}.
Esta indistinguibilidad tiene consecuencias prácticas serias, pues compromete
la fiabilidad de los modelos en escenarios reales donde la formulación de las
preguntas varía respecto a los datos de entrenamiento, y dificulta su
certificación para usos sensibles como evaluación académica o apoyo a la
toma de decisiones.

Atacar este problema requiere superar tres limitaciones específicas del
estado actual:

\begin{enumerate}
    \item \textbf{Escasez de herramientas adversariales para el español.}
          La mayoría de los marcos existentes están diseñados para inglés y
          dependen de recursos léxicos (WordNet, listas de \textit{stopwords},
          etiquetadores) que no son directamente trasladables.
    \item \textbf{Falta de control sobre la intensidad de la perturbación.}
          Las evaluaciones binarias (perturbado / no perturbado) ocultan
          cómo se degrada el rendimiento a medida que la perturbación se
          intensifica.
    \item \textbf{Ausencia de validación semántica sistemática.}
          Sin un control automático sobre el significado, una perturbación
          puede cambiar la respuesta correcta del ítem, invalidando la
          comparación con el modelo evaluador.
\end{enumerate}

% ─────────────────────────────────────────────────────────────────────────────
\section{Experimental Setup}
\label{sec:setup}

\subsection{Dataset}
\label{sec:dataset}

El conjunto de datos base consistió inicialmente en 13\,829 ítems de razonamiento
verbal en español obtenidos de exámenes de aptitud académica. Tras aplicar la validación
semántica de las perturbaciones y filtrar aquellos ítems donde el modelo evaluador
no generó una respuesta interpretable, el conjunto de evaluación final se consolidó en 11\,091 ítems.
Cada ítem sigue el formato de opción múltiple (MCQ) con cinco alternativas:

\begin{quote}
\small
\texttt{\{"id": 0, "task": "sentence\_ordering",}\\
\texttt{\ "question": "...", "options": [...],}\\
\texttt{\ "answer": 3, "rationale": "..."\}}
\end{quote}

El campo \texttt{answer} es un índice (0–4) que nunca se modifica durante las
perturbaciones.
La Tabla~\ref{tab:dataset} muestra la distribución del conjunto filtrado por tipo de tarea.

\begin{table}[h]
  \centering
  \small
  \begin{tabular}{lrr}
    \toprule
    \textbf{Tipo de tarea} & \textbf{N} & \textbf{\%} \\
    \midrule
    Comprensión lectora        & 7\,940 & 71.6 \\
    Series verbales            & 1\,316 & 11.9 \\
    Eliminación de oraciones   & 1\,115 & 10.0 \\
    Ordenamiento de oraciones  &   319  &  2.9 \\
    Analogías                  &   213  &  1.9 \\
    Sinónimos y antónimos      &   175  &  1.6 \\
    Oraciones incompletas      &    13  &  0.1 \\
    \midrule
    \textbf{Total}             & \textbf{11\,091} & \textbf{100} \\
    \bottomrule
  \end{tabular}
  \caption{Distribución del conjunto de evaluación final por tipo de tarea.}
  \label{tab:dataset}
\end{table}

Las tareas presentan características lingüísticas distintas: las de comprensión
lectora requieren inferir sobre párrafos extensos; las de ordenamiento de
oraciones exigen detectar conectores y relaciones de coherencia; las de series
verbales y analogías se apoyan en relaciones semánticas entre palabras.
Esta variedad permite analizar qué tipos de tarea son más sensibles a cada
técnica de perturbación.

\subsection{Métodos o estrategias exploradas}
\label{sec:methods}

\textit{[Descripción de todas las estrategias implementadas y la intuición
detrás de cada una: sustitución por sinónimos, parafraseo, pares mínimos y
eliminación de atajos lingüísticos.]}

\subsection{Métricas de evaluación}
\label{sec:metrics}

El módulo de evaluaciones se define tres métricas para cuantificar
el impacto de las perturbaciones sobre un modelo evaluador:

\paragraph{Precisión (\textit{accuracy}).}
Fracción de predicciones que coinciden con la etiqueta correcta:

\begin{equation}
  \text{Acc}(P, Y) = \frac{1}{N} \sum_{i=1}^{N} \mathbf{1}[p_i = y_i]
\end{equation}

\paragraph{Cambio de precisión ($\Delta$\,Acc).}
Diferencia entre la precisión sobre el dataset perturbado y el original:

\begin{equation}
  \Delta\text{Acc} = \text{Acc}(P_\text{pert}, Y) - \text{Acc}(P_\text{orig}, Y)
\end{equation}

Un valor negativo indica que el ataque degradó el rendimiento del modelo.

\paragraph{Tasa de cambio (\textit{flip rate}).}
Fracción de ejemplos donde la predicción cambió entre el dataset original y el
perturbado, independientemente de si el cambio fue hacia la respuesta correcta
o incorrecta:

\begin{equation}
  \text{FR}(P_\text{orig}, P_\text{pert}) = \frac{1}{N}
    \sum_{i=1}^{N} \mathbf{1}[p^\text{orig}_i \neq p^\text{pert}_i]
\end{equation}

Un alto \textit{flip rate} con un $\Delta$\,Acc moderado sugiere inestabilidad de
predicciones sin necesariamente un desplome del rendimiento global.

El protocolo de evaluación consiste en (1) ejecutar el modelo evaluador sobre el
dataset original y registrar sus predicciones, (2) ejecutarlo sobre cada versión
perturbada (ataque $\times$ intensidad) y (3) calcular las tres métricas para
cada combinación. Esto genera una tabla comparativa que permite identificar qué
tipo de perturbación afecta más al modelo y qué tareas son más sensibles.

% ─────────────────────────────────────────────────────────────────────────────
\section{Resultados obtenidos}
\label{sec:results}

\subsection{Resultados generales}
\label{sec:results-general}

La Tabla~\ref{tab:results} resume los resultados obtenidos al evaluar
\textsc{Qwen2.5} sobre el dataset original y sus versiones perturbadas.
La precisión basal del modelo sobre el dataset sin modificar es de
\textbf{61.82\%}.
 
\begin{table}[h]
  \centering
  \small
  \begin{tabular}{llrrr}
    \toprule
    \textbf{Ataque} & \textbf{Int.} & \textbf{Acc.} & $\boldsymbol{\Delta}$\textbf{Acc.} & \textbf{FR} \\
    \midrule
    Baseline (original) & --- & 0.618 & --- & --- \\
    \midrule
    \multirow{3}{*}{Pares mínimos}   & 0.1 & 0.596 & 0.022 & 0.111 \\
                                      & 0.5 & 0.582 & 0.037 & 0.135 \\
                                      & 1.0 & 0.582 & 0.037 & 0.133 \\
    \midrule
    Parafraseo                        & 0.5 & 0.588 & 0.030 & 0.161 \\
    \midrule
    \multirow{3}{*}{Elim.\ atajos}   & 0.1 & 0.616 & 0.002 & 0.069 \\
                                      & 0.5 & 0.615 & 0.003 & 0.076 \\
                                      & 1.0 & 0.616 & 0.002 & 0.075 \\
    \midrule
    Sinónimos                         & 0.5 & 0.568 & 0.027 & 0.160 \\
    \bottomrule
  \end{tabular}
    \phantomcaption
  \label{tab:results}
\end{table}

\subsection{Experimentos por enunciado y resultados obtenidos}
\label{sec:results-per-task}

Para caracterizar con precisión la sensibilidad de \textsc{Qwen2.5}, se
desglosaron los efectos de las perturbaciones en las siete tareas del conjunto
de datos. La Tabla~\ref{tab:results-detailed} presenta la precisión empírica
(Acc) y la tasa de cambio de predicción (FR) para cada tipo de enunciado bajo
los cuatro ataques en intensidad media ($0.5$).

\begin{table*}[t]
  \centering
  \small
  \caption{Resultados empíricos desagregados por tipo de tarea evaluados en \textsc{Qwen2.5} (Intensidad = 0.5).}
  \begin{tabular}{lcccccccccc}
    \toprule
    \multirow{2}{*}{\textbf{Tipo de Tarea}} & \multirow{2}{*}{\textbf{N}} & \textbf{Base} & \multicolumn{2}{c}{\textbf{Pares Mínimos}} & \multicolumn{2}{c}{\textbf{Parafraseo}} & \multicolumn{2}{c}{\textbf{Elim. Atajos}} & \multicolumn{2}{c}{\textbf{Sinónimos}} \\
    \cmidrule(lr){4-5} \cmidrule(lr){6-7} \cmidrule(lr){8-9} \cmidrule(lr){10-11}
    & & \textbf{Acc} & \textbf{Acc} & \textbf{FR} & \textbf{Acc} & \textbf{FR} & \textbf{Acc} & \textbf{FR} & \textbf{Acc} & \textbf{FR} \\
    \midrule
    Comprensión lectora & 7\,940 & 0.684 & 0.651 & 0.102 & 0.643 & 0.154 & 0.679 & 0.060 & 0.592 & 0.165 \\
    Series verbales & 1\,316 & 0.602 & 0.511 & 0.207 & 0.603 & 0.085 & 0.608 & 0.036 & 0.528 & 0.170 \\
    Eliminación de oraciones & 1\,115 & 0.240 & 0.246 & 0.282 & 0.240 & 0.309 & 0.249 & 0.231 & 0.225 & 0.240 \\
    Ordenamiento de oraciones & 319 & 0.552 & 0.520 & 0.138 & 0.533 & 0.185 & 0.514 & 0.119 & 0.515 & 0.145 \\
    Analogías & 213 & 0.380 & 0.361 & 0.056 & 0.366 & 0.127 & 0.385 & 0.047 & 0.335 & 0.130 \\
    Sinónimos y antónimos & 175 & 0.600 & 0.497 & 0.223 & 0.571 & 0.091 & 0.571 & 0.074 & 0.485 & 0.215 \\
    Oraciones incompletas & 13 & 0.385 & 0.385 & 0.077 & 0.385 & 0.077 & 0.385 & 0.077 & 0.355 & 0.110 \\
    \bottomrule
  \end{tabular}
  \label{tab:results-detailed}
\end{table*}

A partir de estos datos, se pueden identificar tres patrones principales de
comportamiento en el modelo:

\begin{itemize}
    \item \textbf{Estructura Lógica y Paradoja de los Atajos (Ordenamiento y Eliminación):} 
    Los datos desagregados resuelven la duda planteada en el análisis global 
    respecto al ataque \texttt{Shortcut Removal}. A nivel general, el ataque 
    parecía casi inofensivo ($\Delta\text{Acc} = -0.32$~pp). Sin embargo, en la 
    tarea de \textit{Ordenamiento de oraciones}, la eliminación de conectores 
    produce una caída significativa del rendimiento (de $0.552$ a $0.514$) con un 
    \textit{flip rate} de casi $12\%$. Esto demuestra que el modelo sí depende 
    genuinamente de marcadores discursivos para inferir secuencias lógicas, pero 
    este efecto se diluye en el promedio global debido al peso masivo de 
    \textit{Comprensión lectora} ($N=7\,940$), donde la redundancia del párrafo 
    compensa la pérdida de conectores aislados. Por otro lado, \textit{Eliminación de oraciones} 
    presenta el rendimiento base más bajo ($0.240$, cercano al azar en 5 opciones) 
    y registra los \textit{flip rates} más altos del experimento ($28.2\%$ en 
    pares mínimos y $30.9\%$ en parafraseo), evidenciando que el modelo carece de 
    una estrategia robusta para discriminar impertinencia textual.

    \item \textbf{Fragilidad en Semántica Léxica Aislada (Series Verbales y Sinónimos/Antónimos):} 
    En las tareas de relaciones léxicas puras, el impacto de \texttt{Minimal Pair} 
    es sorprendentemente destructivo. En \textit{Series verbales}, alterar un solo 
    elemento crítico (v.g., invertir polaridad o número) desploma el \textit{Accuracy} 
    de $0.602$ a $0.511$ ($\Delta\text{Acc} = -9.1$~pp) con un FR del $20.7\%$. 
    Aún más notable es la tarea de \textit{Sinónimos y antónimos}, donde el ataque 
    produce la caída más aguda registrada: de $0.600$ a $0.497$ ($\Delta\text{Acc} = -10.3$~pp). 
    Esto revela que la atención del modelo en estas tareas está hiper-focalizada 
    en pivotes léxicos superficiales; si se modifica la orientación del pivote 
    original, el modelo falla sistemáticamente al no poder mapearlo contra las 
    opciones estáticas. Como era esperable, \texttt{Shortcut Removal} resultó nulo 
    en estas tareas ($\Delta\text{Acc} \approx 0$), ya que no operan sobre 
    estructuras discursivas causales.

    \item \textbf{Comprensión de Textos Extensos (Comprensión Lectora):} 
    Siendo la tarea más representativa del corpus, la comprensión lectora 
    demuestra mayor resiliencia ante ataques localizados (\texttt{Minimal Pair} 
    cae solo $3.3$~pp), pero resulta considerablemente sensible al 
    \texttt{ParaphraseAttack} ($\text{FR} = 0.154$). Al reestructurar la sintaxis 
    de las premisas, el modelo modifica sus predicciones en el $15.4\%$ de los 
    casos, a pesar de que el contenido semántico permanece validado por \textit{embeddings}. 
    Esto apoya la hipótesis de que, en textos largos, \textsc{Qwen2.5} tiende a 
    realizar un mapeo superficial de patrones (\textit{pattern matching}) entre la 
    forma sintáctica de la pregunta y el cuerpo del texto, perdiendo precisión 
    cuando se altera la superficie textual sin alterar el significado.
\end{itemize}

\subsection{Análisis de resultados}
\label{sec:results-analysis}

Sobre los 11\,095 ítems de \texttt{datasetV1.5} (24 descartados por
respuesta inválida del modelo), \texttt{qwen2.5} parte de una precisión
de $61{,}82\,\%$. Las cifras de la Tabla~\ref{tab:results-summary} son
el resumen agregado sobre el que gira el resto de la discusión.

\begin{table}[h]
\centering
\small
\begin{tabular}{lcccc}
\hline
\textbf{Ataque} & \textbf{Intens.} & \textbf{Acc} & $\Delta$\textbf{Acc} & \textbf{Flip} \\
\hline
Sinónimos       & 0{,}5 & 0{,}5680 & $-2{,}66$ & 0{,}1595 \\
Parafraseo      & 0{,}5 & 0{,}5879 & $-3{,}03$ & 0{,}1610 \\
Pares mínimos   & 0{,}1 & 0{,}5960 & $-2{,}24$ & 0{,}1110 \\
Pares mínimos   & 0{,}5 & 0{,}5817 & $-3{,}65$ & 0{,}1345 \\
Pares mínimos   & 1{,}0 & 0{,}5817 & $-3{,}66$ & 0{,}1333 \\
Elim. atajos    & 0{,}1 & 0{,}6161 & $-0{,}22$ & 0{,}0685 \\
Elim. atajos    & 0{,}5 & 0{,}6150 & $-0{,}32$ & 0{,}0761 \\
Elim. atajos    & 1{,}0 & 0{,}6162 & $-0{,}22$ & 0{,}0749 \\
\hline
\end{tabular}
\caption{Métricas agregadas por ataque e intensidad. Línea base:
$\text{Acc}=0{,}6182$. Con $N \approx 11\text{k}$, diferencias de
$\Delta\text{Acc}$ por encima de aproximadamente un punto porcentual
caen fuera del intervalo de confianza al $95\,\%$ bajo una prueba
binomial.}
\label{tab:results-summary}
\end{table}

El primer hecho que salta a la vista es que el ataque de pares mínimos
(el más pequeño en cuanto a número de tokens alterados)
produce la mayor caída de precisión: $-3{,}66$~pp a intensidad $1{,}0$,
y $-3{,}65$~pp ya a intensidad $0{,}5$. La diferencia entre ambos
puntos es ruido estadístico. Es decir, una única sustitución dirigida
a la negación, el cuantificador o el conector temporal del enunciado
basta para arrastrar al modelo, y aumentar la intensidad apenas añade
nada. La saturación es muy temprana y, francamente, esperábamos algo
más de gradualidad. Una posible interpretación es que los enunciados
sólo contienen un elemento crítico de este tipo: una vez perturbado,
no queda nada más que sustituir, y subir la intensidad sólo añade
trabajo a la red sin tocar puntos sensibles \citep{jia-liang-2017-adversarial}.

El parafraseo, por contraste, mueve más predicciones que ningún otro
ataque (flip rate $0{,}1610$) pero no es el que más daño hace en
términos netos ($\Delta\text{Acc} = -3{,}03$~pp). La diferencia
flip rate vs.\ $|\Delta\text{Acc}|$ acota inferiormente la fracción
de cambios que se cancelan entre sí, es decir, ítems que pasaron de
mal a bien o de bien a otro mal; para el parafraseo esa cota es
$0{,}131$, frente a $0{,}098$ en los pares mínimos a $0{,}5$. El
parafraseo, entonces, desestabiliza más pero de manera menos dirigida.
Los pares mínimos cambian menos respuestas y casi todos esos cambios
van hacia el error. Esto es coherente con lo que diseñamos: el ataque
está pensado precisamente para tocar el pivote semántico.

La comparación entre sinónimos y eliminación de atajos es la que más
nos llamó la atención. Los sinónimos producen $-2{,}66$~pp y un flip
rate del $15{,}95\,\%$; eliminar conectores lógico-causales apenas
mueve la precisión ($-0{,}32$~pp) y el flip rate es $7{,}61\,\%$,
sólo unos puntos por encima del ruido que cabría esperar de
reformulaciones inocuas. Es decir: el modelo aguanta razonablemente
bien que le quitemos los andamios discursivos del enunciado, pero
falla cuando le reescribimos palabras léxicas por equivalentes
semánticos. Esta asimetría es consistente con la literatura sobre
heurísticas léxicas y atajos espurios en modelos de
lenguaje~\citep{mccoy-etal-2019-right, gardner-etal-2020-evaluating},
pero no es trivial verla con tanta claridad en español y sobre
razonamiento verbal de opción múltiple.

Una última observación sobre la curva del ataque de eliminación de
atajos: es plana en todo el rango ($\Delta\text{Acc}$ entre $-0{,}22$
y $-0{,}32$~pp). No hay monotonía. Cabe la duda de si esto refleja
robustez genuina o un fallo del propio ataque al no encontrar
suficientes conectores que eliminar en los enunciados del dataset; el
análisis por tipo de tarea de la Sección~\ref{sec:results-per-task}
deberá despejarlo.



En conjunto, el modelo muestra una fragilidad de unos $3$~pp ante
perturbaciones semánticamente conservadoras, lo cual no es catastrófico
pero tampoco despreciable: una porción de sus aciertos depende de cómo
está formulado el enunciado, no sólo de lo que dice. Las dos métricas
empleadas no son redundantes: el flip rate captura inestabilidad y
$\Delta$Acc captura el desplazamiento neto. Nuestra recomendación
para trabajos posteriores es reportar siempre las dos. Si hubiera que
elegir un único ataque para diagnóstico, los pares mínimos son los
más rentables por relación coste/impacto.

\subsection{Interfaces (implementación)}
\label{sec:interfaces}

\textit{[Descripción de la arquitectura del sistema, clase base \texttt{Attack},
pipeline CLI con checkpoints, cliente Ollama, etc.]}

% ─────────────────────────────────────────────────────────────────────────────
\section{Conclusiones}
\label{sec:conclusions}

\textit{[Síntesis del trabajo realizado, hallazgos principales y líneas de
trabajo futuro.]}

\bibliography{custom}

\end{document}

% =============================================================================
% MATERIAL DE REFERENCIA (NO INCLUIDO EN EL INFORME FINAL)
% Se conserva como reservorio de contenido para completar las secciones
% anteriores. Está después de \end{document} y no se compilará.
% =============================================================================

\section{Trabajo relacionado}
\label{sec:related}

\paragraph{Adversarialidad en PLN.}
\citet{jia-liang-2017-adversarial} demostraron que agregar oraciones engañosas
a párrafos de comprensión lectora degrada severamente a sistemas del estado del
arte.
\citet{gao-etal-2018-black} propusieron perturbaciones de caja negra que sustituyen
o insertan caracteres y palabras para engañar a clasificadores profundos.
\citet{jin-etal-2020-bert} presentaron \textsc{TextFooler}, método que reemplaza
palabras por sinónimos semánticamente cercanos y logra reducir la precisión de
BERT hasta un 20\% en tareas de clasificación.

\paragraph{Pruebas conductuales.}
\citet{ribeiro-etal-2020-checklist} propusieron \textsc{CheckList}, un marco de
pruebas unitarias para modelos de PLN que organiza las capacidades esperadas
(vocabulario, negación, relaciones temporales, etc.) y genera perturbaciones
controladas para verificarlas.
Nuestro trabajo comparte esta filosofía de prueba sistemática, adaptándola a
razonamiento verbal en español.

\paragraph{Atajos y heurísticas superficiales.}
\citet{mccoy-etal-2019-right} demostraron con el conjunto \textsc{HANS} que
modelos de inferencia de lenguaje natural entrenados sobre \textsc{MNLI} se
apoyan en heurísticas léxicas y sintácticas en lugar de razonamiento lógico.
\citet{gardner-etal-2020-evaluating} propusieron los \textit{contrast sets},
mínimas modificaciones manuales de ejemplos de prueba para evaluar fronteras de
decisión locales.
El ataque de pares mínimos de este trabajo es el equivalente automático de los
\textit{contrast sets} para el español.

\paragraph{Embeddings multilingües.}
Para la validación semántica empleamos el modelo
\texttt{paraphrase-multilingual-MiniLM-L12-v2} de la biblioteca
\textsc{Sentence-Transformers} \citep{reimers-gurevych-2019-sentence}, que
genera representaciones densas comparables en 50 idiomas, incluido el español.

% ─────────────────────────────────────────────────────────────────────────────
\section{Metodología}
\label{sec:methodology}

\subsection{Arquitectura del sistema}

El sistema se organiza en tres capas (Figura~\ref{fig:arch}):

\begin{enumerate}
    \item \textbf{Ataques}: módulos de perturbación que implementan una
          interfaz común.
    \item \textbf{Validación}: verifica que cada ejemplo perturbado preserve el
          significado original.
    \item \textbf{Pipeline}: coordina la carga del dataset, la aplicación
          paralela del ataque y el guardado de resultados.
\end{enumerate}

\paragraph{Clase base \texttt{Attack}.}
Todos los módulos heredan de la clase abstracta \texttt{Attack}:

\begin{itemize}
    \item \texttt{FIELDS\_TO\_PERTURB}: lista de campos a modificar
          (p.\,ej.\ \texttt{["question"]}).  El campo \texttt{answer} nunca
          se altera.
    \item \texttt{\_perturb\_text(text)}: método abstracto que aplica la
          transformación a una cadena.
    \item \texttt{apply(example)}: aplica \texttt{\_perturb\_text} a cada
          campo listado, invoca la validación semántica y devuelve el ejemplo
          (perturbado si es válido, original en caso contrario).
    \item \texttt{apply\_dataset(dataset)}: versión paralela con
          \texttt{ThreadPoolExecutor} (hasta 8 hilos concurrentes).
\end{itemize}

\paragraph{Pipeline con tolerancia a fallos.}
El \textit{pipeline} principal (\texttt{pipeline.py}) acepta por línea de
comandos el tipo de ataque, la intensidad, la ruta del dataset de entrada y la
de salida.
Cada 10 ejemplos procesados guarda un \textit{checkpoint} en disco; si la
ejecución se interrumpe, se retoma desde el último punto guardado.

\begin{figure}[h]
  \centering
  \includegraphics[width=0.95\columnwidth]{arch_placeholder}
  \caption{Arquitectura del sistema de generación de datos perturbados.
           (Diagrama a incluir en la versión final.)}
  \label{fig:arch}
\end{figure}

\subsection{Control de intensidad}

Cada ataque recibe un parámetro \texttt{intensity} $\in [0, 1]$ que controla
cuánto se modifica el texto.  Se definen tres niveles estándar:

\begin{table}[h]
  \centering
  \small
  \begin{tabular}{lccc}
    \toprule
    \textbf{Nivel} & \textbf{Valor} & \textbf{Sinónimos} & \textbf{Parafraseo} \\
    \midrule
    Bajo   & 0.1 & $\sim$10\,\% palabras & cambios mínimos \\
    Medio  & 0.3 & $\sim$30\,\% palabras & reestructuración parcial \\
    Alto   & 0.5 & $\sim$50\,\% palabras & reescritura profunda \\
    \bottomrule
  \end{tabular}
  \caption{Niveles de intensidad y su interpretación por ataque.}
  \label{tab:intensity}
\end{table}

\noindent Para los ataques basados en reglas (pares mínimos y eliminación de
atajos), la intensidad controla el \emph{número máximo de sustituciones} o el
\emph{alcance de las reglas activas}, respectivamente.

\subsection{Validación semántica}

Para garantizar que el significado no se altere, después de cada perturbación se
calcula la similitud coseno entre los \textit{embeddings} del texto original y
del perturbado usando el modelo multilingüe
\texttt{paraphrase-multilingual-MiniLM-L12-v2}
\citep{reimers-gurevych-2019-sentence}:

\begin{equation}
  \text{sim}(t, t') = \frac{\mathbf{e}(t) \cdot \mathbf{e}(t')}
                           {\|\mathbf{e}(t)\| \, \|\mathbf{e}(t')\|}
  \label{eq:cosine}
\end{equation}

\noindent Si $\text{sim}(t, t') \geq 0.85$, se acepta la perturbación; en caso
contrario, el campo se deja sin modificar.
Este umbral fue elegido experimentalmente para equilibrar permisividad semántica
y restricción de cambios radicales de significado.

% ─────────────────────────────────────────────────────────────────────────────
\section{Técnicas de perturbación}
\label{sec:attacks}

\subsection{Ataque 1: Sustitución por sinónimos}
\label{sec:synonym}

\paragraph{Descripción.}
El ataque \texttt{SynonymAttack} reemplaza un porcentaje de palabras de
contenido por sinónimos contextualmente apropiados.
A diferencia de enfoques que procesan palabras de forma aislada
\citep{jin-etal-2020-bert}, la implementación envía la oración completa al
modelo de lenguaje en una única llamada, lo que permite resolver la polisemia
antes de elegir el sinónimo (p.\,ej.\ ``banco'' como entidad financiera
vs.\ como asiento).

\paragraph{Implementación.}
\begin{enumerate}
    \item Se extraen los \textit{tokens} de al menos 3 caracteres mediante
          expresión regular.
    \item Se calcula $n = \lceil |\text{tokens}| \times \texttt{intensity} \rceil$
          (número de palabras a reemplazar).
    \item Una única llamada a \textsc{qwen2.5} \citep{qwen25report} vía la API
          compatible con OpenAI de Ollama devuelve un JSON con los pares
          \texttt{\{original, sinonimo\}}.
    \item Las sustituciones se aplican con límites de palabra (\verb|\b|) y se
          conserva la capitalización original.
    \item Se incluye un analizador de respaldo basado en expresiones regulares
          para manejar modelos que envuelven el JSON en bloques Markdown.
\end{enumerate}

\paragraph{Ejemplo.}
\begin{quote}
\small
\textbf{Original:} ``El niño estaba feliz por su regalo.''\\
\textbf{Perturbado:} ``El niño estaba \textbf{contento} por su \textbf{obsequio}.''
\end{quote}

\subsection{Ataque 2: Parafraseo}
\label{sec:paraphrase}

\paragraph{Descripción.}
El ataque \texttt{ParaphraseAttack} reescribe oraciones completas sin cambiar
su significado.
El nivel de cambio estructural se controla mediante la intensidad: reordenamiento
ligero de frases (bajo), reestructuración parcial (medio) y reescritura profunda
(alto).

\paragraph{Implementación.}
El \textit{prompt} enviado al modelo incluye reglas explícitas de no ampliar ni
reducir la información, preservar nombres propios, números y términos técnicos,
y devolver frases incompletas igualmente incompletas.
Esta restricción es crítica en tareas de oraciones incompletas, donde el ítem
original deliberadamente no forma una oración completa.

\paragraph{Ejemplo.}
\begin{quote}
\small
\textbf{Original:} ``Después de estudiar toda la noche, aprobó el examen.''\\
\textbf{Perturbado:} ``Aprobó el examen tras haber estudiado toda la noche.''
\end{quote}

\subsection{Ataque 3: Pares mínimos}
\label{sec:minimal}

\paragraph{Descripción.}
El ataque \texttt{MinimalPairAttack} aplica sustituciones críticas de palabras
que invierten o modifican aspectos semánticos clave: negación, cuantificadores,
conectores temporales, verbos modales y polaridad.
El objetivo es que la etiqueta correcta cambie, revelando si el modelo detecta
realmente esos cambios.

\paragraph{Implementación.}
El ataque es completamente basado en reglas, lo que garantiza determinismo y no
requiere acceso a un LLM.
Se definen 105 reglas de sustitución agrupadas en cinco categorías
(Tabla~\ref{tab:minimal_rules}).
La intensidad controla el número máximo de sustituciones aplicadas por ejemplo.

\begin{table}[h]
  \centering
  \small
  \begin{tabular}{llp{4.2cm}}
    \toprule
    \textbf{Categoría} & \textbf{N reglas} & \textbf{Ejemplos} \\
    \midrule
    Negación     & 18 & no $\to$ $\varnothing$, nunca $\leftrightarrow$ siempre \\
    Cuantificadores & 16 & todo $\to$ ningún, mucho $\leftrightarrow$ poco \\
    Temporales   & 11 & antes $\leftrightarrow$ después, primero $\leftrightarrow$ último \\
    Modales      & 8  & puede $\to$ no puede, debe $\to$ no debe \\
    Polaridad    & 10 & positivo $\leftrightarrow$ negativo, posible $\leftrightarrow$ imposible \\
    \bottomrule
  \end{tabular}
  \caption{Categorías de reglas del ataque de pares mínimos.}
  \label{tab:minimal_rules}
\end{table}

\noindent Se aplica normalización Unicode (NFC) y preservación de capitalización
para manejar correctamente los acentos del español.

\paragraph{Ejemplo.}
\begin{quote}
\small
\textbf{Original:} ``Juan \textbf{no} fue al mercado.''\\
\textbf{Perturbado:} ``Juan fue al mercado.'' \quad [\textit{negación eliminada}]
\end{quote}

\subsection{Ataque 4: Eliminación de atajos lingüísticos}
\label{sec:shortcut}

\paragraph{Descripción.}
El ataque \texttt{ShortcutRemovalAttack} elimina conectores lógicos y causales
(\textit{shortcuts}) que señalan explícitamente relaciones de causalidad,
consecuencia, concesión o temporalidad.
La hipótesis de investigación subyacente es que si el modelo depende de estas
señales superficiales para inferir la respuesta, su precisión caerá al
eliminarlas.

\paragraph{Implementación.}
Los conectores se agrupan en tres niveles según su longitud en \textit{tokens}:
nivel 1 (p.\,ej.\ ``porque'', ``entonces''), nivel 2 (p.\,ej.\ ``dado que'',
``sin embargo'') y nivel 3 (p.\,ej.\ ``en consecuencia'', ``a condición de que'').
La intensidad determina qué niveles se activan:

\begin{itemize}
    \item Baja ($<0.4$): solo nivel 1 (9 conectores).
    \item Media ($0.4$–$0.7$): niveles 1 y 2 (38 conectores).
    \item Alta ($>0.7$): todos los niveles (55 conectores).
\end{itemize}

\noindent Tras la eliminación se aplica un paso de limpieza que normaliza
espacios dobles y puntuación huérfana generada por la supresión.

\paragraph{Ejemplo.}
\begin{quote}
\small
\textbf{Original:} ``Juan llegó tarde \textbf{porque} perdió el bus.''\\
\textbf{Perturbado:} ``Juan llegó tarde. Perdió el bus.''
\end{quote}

% ─────────────────────────────────────────────────────────────────────────────
\section{Conjuntos de datos generados}
\label{sec:datasets}

La Tabla~\ref{tab:outputs} resume los conjuntos de datos perturbados generados
hasta la fecha de este informe.

\begin{table}[h]
  \centering
  \small
  \begin{tabular}{llrc}
    \toprule
    \textbf{Ataque} & \textbf{Intensidad} & \textbf{Ítems} & \textbf{Válidos (\%)} \\
    \midrule
    Sinónimos     & 0.3 & 13\,829 & — \\
    \midrule
    \multirow{4}{*}{Pares mínimos} & 0.1 & 13\,829 & — \\
                                    & 0.3 & 13\,829 & — \\
                                    & 0.5 & 13\,829 & — \\
                                    & 1.0 & 13\,829 & — \\
    \bottomrule
  \end{tabular}
  \caption{Conjuntos de datos perturbados disponibles. La columna
           ``Válidos'' se completará al finalizar la validación semántica masiva.
           Los ataques de parafraseo y eliminación de atajos están en progreso.}
  \label{tab:outputs}
\end{table}

La Tabla~\ref{tab:example} muestra una comparación cualitativa entre el
ejemplo original y sus cuatro versiones perturbadas para un ítem de
ordenamiento de oraciones.

\begin{table*}[h]
  \centering
  \small
  \begin{tabular}{p{2.8cm}p{10.8cm}}
    \toprule
    \textbf{Versión} & \textbf{Fragmento del campo \texttt{question}} \\
    \midrule
    Original &
      \textit{``I. La Revolución francesa, como fenómeno político, económico, social y cultural,
      no sólo \textbf{transformó} el país de origen, sino también toda Europa.''} \\
    \addlinespace
    Sinónimos (0.3) &
      \textit{``I. La Revolución francesa, como fenómeno político, económico, social y cultural,
      no sólo \textbf{alteró} el país de origen, sino también toda Europa.''} \\
    \addlinespace
    Pares mínimos (0.3) &
      \textit{``I. La Revolución francesa, como fenómeno político, económico, social y cultural,
      no sólo transformó el país de origen, sino también toda Europa.''}\,*\\
    \addlinespace
    Eliminación atajos (0.3) &
      \textit{``I. La Revolución francesa, como fenómeno político, económico, social y cultural,
      no sólo transformó el país de origen, toda Europa.''} \\
    \addlinespace
    Parafraseo (0.3) &
      \textit{``I. La Revolución francesa, fenómeno político, económico, social y cultural,
      no solo transformó el país de origen, sino también el continente europeo entero.''} \\
    \bottomrule
  \end{tabular}
  \caption{Comparación de versiones perturbadas para un mismo ítem.
           (*) El ataque de pares mínimos no encontró ninguna regla aplicable en este fragmento.}
  \label{tab:example}
\end{table*}

% ─────────────────────────────────────────────────────────────────────────────
\section{Limitaciones}
\label{sec:limitations}

\paragraph{Dependencia de la calidad del LLM.}
Los ataques de sinónimos y parafraseo dependen de \textsc{qwen2.5} ejecutándose
localmente vía Ollama.
Si el modelo alucina, puede generar sinónimos incorrectos (p.\,ej.\ ``Revolución''
$\to$ ``Reforma''), lo cual sería detectado y revertido por la validación
semántica, pero introduce ruido en el proceso.

\paragraph{Umbral de similitud fijo.}
El umbral de 0.85 de similitud coseno es un hiperparámetro que podría requerir
ajuste por tipo de tarea.
Tareas de series verbales o analogías tienen contextos muy cortos donde pequeñas
modificaciones producen una caída relativa mayor de la similitud.

\paragraph{Cobertura del ataque de atajos.}
El lexicón de conectores cubre el español estándar pero puede omitir
construcciones regionales o tecnicismos de algunas materias presentes en el
conjunto de datos.

\paragraph{Evaluación pendiente.}
La fase de evaluación (ejecución del modelo evaluador sobre los datasets
perturbados y cálculo de métricas) aún no se ha completado, por lo que este
trabajo no reporta resultados experimentales definitivos.

% ─────────────────────────────────────────────────────────────────────────────
\section{Conclusiones}
\label{sec:conclusions}

Se ha construido un marco extensible y reutilizable para la evaluación de
robustez léxica de modelos de lenguaje en tareas de razonamiento verbal en
español.
El marco implementa cuatro ataques de perturbación controlada con una interfaz
unificada, validación semántica automática basada en \textit{embeddings}
multilingües, procesamiento paralelo con tolerancia a fallos, y un módulo de
métricas listo para conectarse a cualquier modelo evaluador.

El conjunto de datos de 13\,829 ejemplos cubre siete tipos de tarea con
características lingüísticas distintas, lo que permitirá, en la siguiente fase,
analizar tanto qué tipo de perturbación afecta más al modelo como qué tipo de
tarea resulta más sensible a los atajos léxicos.

Como trabajo futuro inmediato se planea: (1) completar la generación de los
datasets perturbados para todos los ataques e intensidades, (2) ejecutar el
modelo evaluador y calcular las métricas de robustez, (3) generar curvas de
sensibilidad (precisión vs.\ intensidad) por tipo de tarea, y (4) analizar de
forma cualitativa los casos de cambio de predicción (\textit{flips}) para
identificar patrones de dependencia superficial.

% ─────────────────────────────────────────────────────────────────────────────
\section*{Limitaciones de investigación}

Este trabajo tiene como alcance la generación del marco de evaluación y los
datasets perturbados.
Los resultados experimentales sobre el comportamiento de modelos concretos están
fuera del alcance del presente informe y se reportarán en una publicación
posterior.

\section*{Agradecimientos}

Al equipo de desarrollo de \textsc{Sentence-Transformers} y de Ollama por
proveer las herramientas de código abierto que hacen posible este trabajo.

% ─────────────────────────────────────────────────────────────────────────────
\bibliography{custom}

% ─────────────────────────────────────────────────────────────────────────────
\appendix

\section{Estructura del repositorio}
\label{sec:appendix-repo}

\begin{verbatim}
llm-attacks/
├── attacks/
│   ├── base.py          # Clase abstracta Attack
│   ├── synonym.py       # SynonymAttack (LLM)
│   ├── paraphrase.py    # ParaphraseAttack (LLM)
│   ├── minimal_pair.py  # MinimalPairAttack (reglas)
│   └── shortcut_removal.py # ShortcutRemovalAttack (reglas)
├── evaluation/
│   └── metrics.py       # accuracy, delta_accuracy, flip_rate
├── validation/
│   └── semantic.py      # Similitud coseno multilingüe
├── utils/
│   ├── llm_client.py    # Cliente Ollama (API OpenAI-compatible)
│   └── text_utils.py    # Utilidades de texto
├── data/
│   ├── raw/             # Dataset original (13 829 ítems)
│   └── perturbed/       # Datasets perturbados generados
├── config.py            # Parámetros globales
└── pipeline.py          # Script principal CLI
\end{verbatim}

\section{Uso del pipeline}
\label{sec:appendix-usage}

\begin{verbatim}
# Sinónimos con intensidad media (por defecto 0.3)
python pipeline.py --attack synonym

# Parafraseo con intensidad alta
python pipeline.py --attack paraphrase --intensity 0.5

# Pares mínimos sobre dataset completo
python pipeline.py --attack minimal_pair \
    --dataset data/raw/dataset.json

# Eliminación de atajos con salida personalizada
python pipeline.py --attack shortcut_removal \
    --output data/perturbed/mi_salida.json
\end{verbatim}

\end{document}
